\documentclass[a4paper,14pt, unknownkeysallowed]{extreport}
\usepackage{placeins}
\usepackage{cmap} % Улучшенный поиск русских слов в полученном pdf-файле
\usepackage[T2A]{fontenc} % Поддержка русских букв
\usepackage[utf8]{inputenc} % Кодировка utf8
\usepackage[english,russian]{babel} % Языки: русский, английский
%\usepackage{pscyr} % Нормальные шрифты
\usepackage{enumitem}
\usepackage{multirow}
\usepackage[14pt]{extsizes}

\usepackage{caption}
\captionsetup{labelsep=endash}
\captionsetup[figure]{name={Рисунок}}

\usepackage{amsmath}

\usepackage{geometry}
\geometry{left=30mm}
\geometry{right=15mm}
\geometry{top=20mm}
\geometry{bottom=20mm}

\setlength{\parindent}{1.25cm} % для шрифта 14pt

\usepackage{titlesec}
\titleformat{\section}
	{\normalsize\bfseries}
	{\thesection}
	{1em}{}
\titlespacing*{\chapter}{0pt}{-30pt}{8pt}
\titlespacing*{\section}{\parindent}{*4}{*4}
\titlespacing*{\subsection}{\parindent}{*4}{*4}

\usepackage{setspace}
\onehalfspacing % Полуторный интервал

\frenchspacing
\usepackage{indentfirst} % Красная строка

\usepackage{titlesec}
\titleformat{\chapter}{\LARGE\bfseries}{\thechapter}{20pt}{\LARGE\bfseries}
\titleformat{\section}{\Large\bfseries}{\thesection}{20pt}{\Large\bfseries}

\usepackage{listings}
\usepackage{xcolor}

% Для листинга кода:
\lstset{ %
	language=lisp,   					% выбор языка для подсветки	
	basicstyle=\small\sffamily,			% размер и начертание шрифта для подсветки кода
	numbers=left,						% где поставить нумерацию строк (слева\справа)
	%numberstyle=,					% размер шрифта для номеров строк
	stepnumber=1,						% размер шага между двумя номерами строк
	numbersep=5pt,						% как далеко отстоят номера строк от подсвечиваемого кода
	frame=single,						% рисовать рамку вокруг кода
	tabsize=4,							% размер табуляции по умолчанию равен 4 пробелам
	captionpos=t,						% позиция заголовка вверху [t] или внизу [b]
	breaklines=true,					
	breakatwhitespace=true,				% переносить строки только если есть пробел
	escapeinside={\#*}{*)},				% если нужно добавить комментарии в коде
	backgroundcolor=\color{white},
}
\usepackage{pgfplots}
\usetikzlibrary{datavisualization}
\usetikzlibrary{datavisualization.formats.functions}

\usepackage{graphicx}
\newcommand{\img}[3] {
	\begin{figure}[h!]
		\center{\includegraphics[height=#1]{inc/img/#2}}
		\caption{#3}
		\label{img:#2}
	\end{figure}
}
\newcommand{\boximg}[3] {
	\begin{figure}[h]
		\center{\fbox{\includegraphics[height=#1]{inc/img/#2}}}
		\caption{#3}
		\label{img:#2}
	\end{figure}
}

% для матриц
\makeatletter
\renewcommand*\env@matrix[1][\arraystretch]{%
	\edef\arraystretch{#1}%
	\hskip -\arraycolsep
	\let\@ifnextchar\new@ifnextchar
	\array{*\c@MaxMatrixCols c}}
\makeatother

\usepackage{amsfonts}
\usepackage{graphicx}


\usepackage[justification=centering]{caption} % Настройка подписей float объектов

\usepackage[unicode,pdftex]{hyperref} % Ссылки в pdf
\hypersetup{hidelinks}

\usepackage{csvsimple}

\addto\captionsrussian{\def\refname{Список используемых источников}}

\newcommand{\code}[1]{\texttt{#1}}

\makeatletter
\renewcommand{\@biblabel}[1]{#1. \hfill}
\makeatother


\begin{document}


\begin{titlepage}
	\newgeometry{pdftex, left=2cm, right=2cm, top=2.5cm, bottom=2.5cm}
	\fontsize{12pt}{12pt}\selectfont
	\noindent \begin{minipage}{0.15\textwidth}
		\includegraphics[width=\linewidth]{inc/img/b_logo.jpg}
	\end{minipage}
	\noindent\begin{minipage}{0.9\textwidth}\centering
		\textbf{Министерство науки и высшего образования Российской Федерации}\\
		\textbf{Федеральное государственное бюджетное образовательное учреждение высшего образования}\\
		\textbf{«Московский государственный технический университет имени Н.Э.~Баумана}\\
		\textbf{(национальный исследовательский университет)»}\\
		\textbf{(МГТУ им. Н.Э.~Баумана)}
	\end{minipage}
	
	\noindent\rule{18cm}{3pt}
	\newline\newline
	\noindent ФАКУЛЬТЕТ $\underline{\text{«Информатика и системы управления»~~~~~~~~~~~~~~~~~~~~~~~~~~~~~~~~~~~~~~~~~~~~~~~~~~~~~~~}}$ \newline\newline
	\noindent КАФЕДРА $\underline{\text{«Программное обеспечение ЭВМ и информационные технологии»~~~~~~~~~~~~~~~~~~~~~~~}}$\newline\newline\newline\newline\newline\newline\newline
	
	
	\begin{center}
		\Large\textbf{Отчет по лабораторной работе №2}
		\Large\textbf {по дисциплине <<Функциональное и логическое программирование>>}
	\end{center}
	
	\noindent\textbf{Тема} $\underline{\text{~~Определений функций пользователя~~~~~~~~~~~~~~~~~~~~~~~~~~~~~~~~~~~~~~~~~~~~~~~~~~~~~~~~~~~~~~~~~}}$\newline\newline\newline
	\noindent\textbf{Студент} $\underline{\text{~~Котляров Н.А.~~~~~~~~~~~~~~~~~~~~~~~~~~~~~~~~~~~~~~~~~~~~~~~~~~~~~~~~~~~~~~~~~~~~~~~~~~~~~~~~~~~~~~~~~}}$\newline\newline
	\noindent\textbf{Группа} $\underline{\text{~~ИУ7-61Б~~~~~~~~~~~~~~~~~~~~~~~~~~~~~~~~~~~~~~~~~~~~~~~~~~~~~~~~~~~~~~~~~~~~~~~~~~~~~~~~~~~~~~~~~~~~~~~~~}}$\newline\newline
	\noindent\textbf{Оценка (баллы)} $\underline{\text{~~~~~~~~~~~~~~~~~~~~~~~~~~~~~~~~~~~~~~~~~~~~~~~~~~~~~~~~~~~~~~~~~~~~~~~~~~~~~~~~~~~~~~~~~~~~~~~~~~~}}$\newline\newline
	\noindent\textbf{Преподаватель} $\underline{\text{~~Толпинская Н.~Б.~~Строганов~Ю.~В.~~~~~~~~~~~~~~~~~~~~~~~~~~~~~~~~~~~~~~~~~~~~~~~~~~}}$\newline
	
	\begin{center}
		\vfill
		Москва~---~\the\year
		~г.
	\end{center}
	\restoregeometry
\end{titlepage}



\section*{Задание 1}

\textbf{Составить диаграмму вычисления следующих выражений::}
\begin{enumerate}[label=\arabic*)]
	\item \lstinline {(equal 3 (abs - 3))}
	\item \lstinline {(equal (+ 1 2) 3)}
	\item \lstinline {(equal (* 4 7) 21)}
	\item \lstinline {(equal (* 2 3) (+ 7 2))}
	\item \lstinline {(equal (- 7 3) (* 3 2))}
	\item \lstinline {(equal (abs (- 2 4)) 3)}
	
\end{enumerate}

\clearpage

\section*{Задание 2}

\textbf{Написать функцию, вычисляющую гипотенузу прямоугольного}

\begin{lstlisting}
(defun calc_hyp(a b)
	(sqrt
		(+ (* a a) (* b b))
	)
)
\end{lstlisting}

\section*{Задание 3}

\textbf{Каковы результаты вычисления следующих выражений?
	(объяснить возможную ошибку и варианты ее устранен6ия)}

\begin{enumerate}[label=\arabic*)]
	\item \lstinline {(list 'a c)} --- 
	Результат: Ошибка. 
	Вариант решения: добавить ко второму аргументу quote.
	\item \lstinline {(cons 'a (b c))} --- 
	Результат: Ошибка. 
	Вариант решения: добавить ко второму аргументу quote.
	\item \lstinline {(cons 'a '(b c))} --- 
	Результат: (a b c) 
	\item \lstinline {(caddr (1 2 3 4 5))} --- 
	Результат: Ошибка. 
	Вариант решения: добавить quote.
	\item \lstinline {(cons 'a 'b 'c)} --- 
	Результат: Ошибка. 
	Вариант решения: убрать третий аргумент.
	\item \lstinline {(list 'a (b c))} --- 
	Результат: Ошибка. 
	Вариант решения: добавить quote к последниму аргументу.
	\item \lstinline {(list a '(b c))}  --- 
	Результат: Ошибка. 
	Вариант решения: добавить quote к первому аргументу.
	\item \lstinline {(list (+ 1 '(length '(1 2 3))))}  --- 
	Результат: Ошибка. 
	Вариант решения: убрать quote перед length. Тогда length будет интепритирована как функция, а не данные.
	
\end{enumerate}

\section*{Задание 4}

\textit{Написать функцию longer\_then от двух списков-аргументов, 
	которая возвращает T, если первый аргумент имеет большую длину.}


\begin{lstlisting}
	(defun longer_then(f s) (> (length f) (length s)))
\end{lstlisting}


\section*{Задание 5}

\textbf{Каковы результаты вычисления следующих выражений?}

\begin{enumerate}[label=\arabic*)]
	\item \lstinline {(cons 3 (list 5 6))} --- (3 5 6)
	\item \lstinline {(list 3 'from 9 'lives (- 9 3))} --- (3 from 9 lives 6)
	\item \lstinline {(+ (length '(for 2 too)) (car '(21 22 23)))} --- 4
	\item \lstinline {(cdr '(cons is short for ans))} --- (is short for ans)
	\item \lstinline {(car (list one two))} --- (Переменна one не определена)
	\item \lstinline {(cons 3 '(list 5 6))} --- (3 list 5 6)
	\item \lstinline {(car (list 'one 'two))} --- one
\end{enumerate}

\section*{Задание 6}

\textbf{Дана функция (defun mystery (x) (list (second x) (first x))).
	Какие результаты вычисления следующих выражений? }


\begin{enumerate}[label=\arabic*)]
	\item \lstinline {(mystery (one two))} --- one не определена.
	\item \lstinline {(mystery (last one two))} --- last не определена.
	\item \lstinline {(mystery free)} --- free не определена.
	\item \lstinline {(mystery one 'two))} --- one не определена.
\end{enumerate}

\section*{Задание 7}

\textbf{Написать функцию, которая переводит температуру в системе Фаренгейта
	температуру по Цельсию (defum f-to-c (temp)…).
	Формулы: 	c = 5/9*(f-32.0); 	f= 9/5*c+32.0. 
	Как бы назывался роман Р.Брэдбери "+451 по Фаренгейту" в системе по Цельсию?}

\begin{lstlisting}
	(defun far_to_cel(temp) (* (/ 5 9)(- temp 32.0)))
\end{lstlisting}

Ответ: 232.77779 по Цельсию

\section*{Задание 8}

\textbf{Что получится при вычисления каждого из выражений?}


\begin{enumerate}[label=\arabic*)]
	\item \lstinline {(list 'cons t NIL)} --- (cons t NIL)
	\item \lstinline {(eval (eval (list 'cons t NIL)))} --- Переменная t не определена
	\item \lstinline {(apply #cons "(t NIL))} --- Неподдерживаемый синтаксис: \#cons
	\item \lstinline {(list 'eval NIL)} --- (eval NIL)
	\item \lstinline {(eval (list 'cons t NIL))} --- (T)
	\item \lstinline {(eval NIL)} --- NIL
	\item \lstinline {(eval (list 'eval NIL))} --- NIL
	
	
\end{enumerate}

\end{document}